% Practice-round swatch sheet: candidate treatments side by side at manuscript scale (150 mm width).
\documentclass[10pt,border=2mm]{standalone}
\usepackage[T1]{fontenc}
\usepackage{figurestyle}
\input{molecules.tex}
\input{numbers.tex}
\begin{document}
\begin{tikzpicture}[plate]
\path[use as bounding box] (0,0) rectangle (15,17.2);
\node[panel] at (0,17.15) {Row 1. Molecule glyphs for the same computed $X_0$, scale 0.6 cm per \AA};
% (a) current: labelled discs
\pic[scale=.6] at (1.9,15.3) {mla-ref};
\node[note] at (1.9,13.95) {(a) current: labels inside discs};
% (b) small discs, labels beside in ink, thinner bonds
\begin{scope}
  \tikzset{atomH/.append style={minimum size=2.2mm,line width=.45pt},atomC/.append style={minimum size=3.0mm},atomN/.append style={minimum size=3.0mm},molbond/.append style={line width=.5pt}}
  \renewcommand\molatom[4]{\node[atom#1] at (#3,#4) {};\if#1H\node[font=\fontsize{8}{9}\selectfont,inner sep=.5pt] at ($(#3,#4)!-.42cm!(0,0)$) {$\mathrm{#1}_{#2}$};\else\node[font=\fontsize{8}{9}\selectfont,inner sep=.5pt,anchor=north] at ($(#3,#4)+(0,-.2)$) {$\mathrm{#1}_{#2}$};\fi}
  \pic[scale=.6] at (6.3,15.3) {mla-ref};
\end{scope}
\node[note] at (6.3,13.95) {(b) small discs, ink labels beside};
% (c) flat structural drawing, connectivity only (not a projection)
\begin{scope}[shift={(11.2,15.3)}]
  \tikzset{atomH/.append style={minimum size=2.2mm,line width=.45pt},atomC/.append style={minimum size=3.0mm},atomN/.append style={minimum size=3.0mm}}
  \draw[molbond,line width=.5pt] (-.65,0)--(.65,0) (-.65,0)--(-1.45,0) (-.65,0)--(-1.1,.75) (-.65,0)--(-1.1,-.75) (.65,0)--(1.15,.75) (.65,0)--(1.15,-.75);
  \node[atomC] at (-.65,0) {}; \node[atomN] at (.65,0) {};
  \foreach \x/\y in {-1.45/0,-1.1/.75,-1.1/-.75,1.15/.75,1.15/-.75} \node[atomH] at (\x,\y) {};
  \node[small] at (-.65,-.38) {$\mathrm C_6$}; \node[small] at (.65,-.38) {$\mathrm N_7$};
  \node[small] at (-1.85,0) {$\mathrm H_1$}; \node[small] at (-1.1,1.1) {$\mathrm H_2$}; \node[small] at (-1.1,-1.1) {$\mathrm H_3$};
  \node[small] at (1.15,1.1) {$\mathrm H_4$}; \node[small] at (1.15,-1.1) {$\mathrm H_5$};
\end{scope}
\node[note] at (11.2,13.95) {(c) flat connectivity drawing (no geometry)};

\node[panel] at (0,13.3) {Row 2. The reduced-family cylinder};
\foreach \cx/\kind in {2.4/1,7.5/2,12.6/3} {
\begin{scope}[shift={(\cx,10.9)}]
  \def\Rx{1.5}\def\Ry{.45}\def\Hh{1.0}
  \ifnum\kind=1
    \draw[contour] (-\Rx,-\Hh)--(-\Rx,\Hh) (\Rx,-\Hh)--(\Rx,\Hh);
    \draw[contour] (0,\Hh) ellipse[x radius=\Rx,y radius=\Ry];
    \draw[contour] (-\Rx,-\Hh) arc[start angle=180,end angle=360,x radius=\Rx,y radius=\Ry];
    \draw[guidedash] (-\Rx,-\Hh) arc[start angle=180,end angle=0,x radius=\Rx,y radius=\Ry];
    \foreach \a/\h in {0/\Hh,120/\Hh,240/\Hh,300/-\Hh,60/-\Hh} \fill[PlateInk] ({\Rx*sin(\a)},{\h-\Ry*cos(\a)}) circle[radius=1.6pt];
    \draw[operation,draw=PlateBlue] plot[domain=0:120,samples=30,variable=\a] ({\Rx*sin(\a)},{\Hh-\Ry*cos(\a)});
  \else
    % light: thin contours, one pale tint on the visible wall, hidden rim dashed, small gold markers with leaders
    \path[fill=ColNitrogen!8] (-\Rx,-\Hh) arc[start angle=180,end angle=360,x radius=\Rx,y radius=\Ry] -- (\Rx,\Hh) arc[start angle=0,end angle=-180,x radius=\Rx,y radius=\Ry] -- cycle;
    \draw[line width=.4pt,draw=PlateInk] (-\Rx,-\Hh)--(-\Rx,\Hh) (\Rx,-\Hh)--(\Rx,\Hh);
    \draw[line width=.4pt,draw=PlateInk,fill=white] (0,\Hh) ellipse[x radius=\Rx,y radius=\Ry];
    \draw[line width=.4pt,draw=PlateInk] (-\Rx,-\Hh) arc[start angle=180,end angle=360,x radius=\Rx,y radius=\Ry];
    \draw[line width=.35pt,draw=PlateInk!60,densely dashed] (-\Rx,-\Hh) arc[start angle=180,end angle=0,x radius=\Rx,y radius=\Ry];
    \draw[line width=.7pt,draw=PlateBlue] plot[domain=0:120,samples=30,variable=\a] ({\Rx*sin(\a)},{\Hh-\Ry*cos(\a)});
    \foreach \a/\h in {0/\Hh,120/\Hh,240/\Hh,300/-\Hh,60/-\Hh} \draw[line width=.4pt,fill=ColGold] ({\Rx*sin(\a)},{\h-\Ry*cos(\a)}) circle[radius=1.4pt];
    \draw[guide] ({\Rx*sin(0)},{\Hh-\Ry*cos(0)})--++(.25,-.35) node[anchor=north west,inner sep=1pt,small] {$H$};
    \draw[guide] ({\Rx*sin(120)},{\Hh-\Ry*cos(120)})--++(.3,.25) node[anchor=west,inner sep=1pt,small] {$tH$};
    \ifnum\kind=3
      % sparse hatch band: a normal-thickening patch on the wall
      \path[hatch] ({\Rx*sin(20)-.2},{-.3-\Ry*cos(20)}) rectangle ({\Rx*sin(20)+.2},{.5-\Ry*cos(20)});
      \draw[guide] ({\Rx*sin(20)-.2},{-.3-\Ry*cos(20)}) rectangle ({\Rx*sin(20)+.2},{.5-\Ry*cos(20)});
    \fi
  \fi
\end{scope}}
\node[note,align=center] at (2.4,9.3) {(a) current: .65 pt contours,\\black dots};
\node[note,align=center] at (7.5,9.3) {(b) light: .4 pt, pale tint,\\gold markers, leaders};
\node[note,align=center] at (12.6,9.3) {(c) as (b) plus\\a sparse hatch patch};

\node[panel] at (0,8.7) {Row 3. Region treatments and the base palette};
\foreach \cx/\kind/\lab in {1.6/1/{sparse hatch},4.4/2/{flat tint},7.2/3/{tint + hatch}} {
\begin{scope}[shift={(\cx,6.9)}]
  \ifnum\kind=1 \path[hatch] (-1.1,-.75) rectangle (1.1,.75); \fi
  \ifnum\kind=2 \path[fill=ColGold!35] (-1.1,-.75) rectangle (1.1,.75); \fi
  \ifnum\kind=3 \path[fill=ColGold!25] (-1.1,-.75) rectangle (1.1,.75); \path[hatch] (-1.1,-.75) rectangle (1.1,.75); \fi
  \draw[line width=.45pt] (-1.1,-.75) rectangle (1.1,.75);
  \node[tag] at (0,0) {$\mathcal F$};
  \node[note] at (0,-1.05) {\lab};
\end{scope}}
\foreach \cx/\col/\lab in {9.6/PlateInk/{ink},10.5/white/{H},11.4/ColOxygen/{O},12.3/ColNitrogen/{N},13.2/ColPotassium/{K},14.1/ColRubidium/{Rb}} {
  \draw[line width=.45pt,fill=\col] (\cx,7.0) circle[radius=.33];
  \node[note] at (\cx,6.35) {\lab};
}
\draw[line width=.45pt,fill=ColGold] (9.6,5.6) circle[radius=.33]; \node[note,anchor=west] at (10.05,5.6) {gold accent, fill};
\draw[line width=.8pt,draw=ColGoldStroke] (9.3,4.95)--(10.2,4.95); \node[note,anchor=west] at (10.3,4.95) {gold stroke, darkened centrally};
\draw[line width=.8pt,draw=ColNitrogen] (9.3,4.45)--(10.2,4.45); \node[note,anchor=west] at (10.3,4.45) {blue stroke};
\draw[line width=.8pt,draw=ColOxygen,densely dashed] (9.3,3.95)--(10.2,3.95); \node[note,anchor=west] at (10.3,3.95) {red dashed stroke};

\node[panel] at (0,4.3) {Row 4. Line hierarchy at this scale};
\foreach \y/\w/\lab in {3.7/.35pt/{guide .35 pt},3.2/.45pt/{light contour .45 pt},2.7/.65pt/{contour .65 pt},2.2/.8pt/{operation .8 pt},1.7/1.0pt/{heavy 1.0 pt}} {
  \draw[line width=\w] (0.2,\y)--(3.2,\y); \node[note,anchor=west] at (3.4,\y) {\lab};
}
\node[note,anchor=north west,text width=8cm,align=left] at (6.5,3.3) {labels 9 pt serif, notes 8 pt; the question for the author: which rows (a), (b), (c) read as "textbook ink with gouache colour" at this size, in colour and in grayscale};
\end{tikzpicture}
\end{document}
